\section{GPT-6.2 AstraとSolを発表}
OpenAIはGPT-6.2 AstraとSolを発表した。Astraは難度の高い専門業務や研究、Solは日常業務と複数工程のエージェント処理を主な対象とする。Grokの収集では、ChatGPT Work、Codex、APIでの提供が案内され、Astraのタスク当たり費用はGPT-6 Astraの3分の1未満、Solはより低費用でAstraに近いコーディング性能を目指すと説明された。API価格についても複数の数値が流通したが、利用条件や実効コストは今後の評価対象となる。初期利用者からは品質への批判も見られた。
\dailyxsource{OpenAI公式投稿}{https://x.com/OpenAI/status/2122750124543926515}

\section{Gemini 4 Hydrogenでロボット制御}
Google DeepMindはGemini 4 HydrogenとAntigravityのロボット向けプラグインを組み合わせ、自然言語からロボットアームを操作するデモを公開した。MuJoCoで動作を練習し、実物の仕分けなどへ展開する流れを紹介している。投稿では、日常的な物体移動で大規模なArgonに匹敵し、空間判断ではGemini Robotics ER 2を上回るとの主張もあった。ただし詳細な技術ブログや実環境での失敗率は観測窓内で確認されていない。
\dailyxsource{Google DeepMind公式投稿}{https://x.com/GoogleDeepMind/status/2122767521925230896}

\section{MAI Code 2のCopilot展開に言及}
MicrosoftのMAI Code 2がGitHub Copilotで一般提供されたとの投稿が見られた。コミュニティではGPT-6.1 Solの98\%に相当する性能を40\%の価格で実現したという比較が共有されたが、再現性のある独立評価は限定的である。GrokはGitHub Changelogへの言及も収集したものの、公式情報の確認時刻は観測窓の境界付近に当たる。実際の提供条件や難問への適性は、今後の利用報告を待つ必要がある。
\dailyxsource{コミュニティ投稿}{https://x.com/LowdCode/status/2122810956971643391}

\section{日立の委任証明PoHDが話題}
日立がAIエージェントへの人間の委任を証明する「Proof of Human Delegation（PoHD）」のプロトタイプを発表したとの投稿があった。生体情報から署名時に鍵を再構築し、鍵を保存しない方式として紹介されている。調達業務を想定し、エージェントによる契約や金銭取引に真正な承認を結び付ける狙いがあるという。Grok原文では日立のプレスリリースに言及するがURLは未確認で、実装詳細は投稿上の説明として扱う。
\dailyxsource{関連投稿}{https://x.com/2logics/status/2122811819912024232}

\section{新モデルの価格と使用感を比較}
GPT-6.2 AstraをClaude Fable 5.2、Gemini 4 Carbon、Grok 4.7と比べる投稿が相次いだ。入力・出力100万トークン当たりの価格比較や、コーディング時の早期終了、応答の省略を批判する声などが共有された。一方、こうした使用感は個別の環境やタスクに左右される。Grokはこれらを未検証と判定しており、価格表やベンチマークを公式の確定値として採用しない。比較条件の統一が今後の論点となる。
\dailyxsource{個人比較の投稿}{https://x.com/tirupathipodugu/status/2122811975612911654}

\section{AIサイバーセキュリティETFの報告}
BitwiseがAIサイバーセキュリティ企業を対象とするETF「ZDAY」をNYSEに上場したとの投稿が観測された。経費率0.75\%、Palo Alto Networksなどを組み入れるとされ、同社初の非暗号資産ETFとの説明もあった。ただしGrokは観測窓内で公式の上場資料を確認していない。AIに伴う防御需要を投資テーマとする動きとして注目されたが、銘柄構成や商品条件は未確認の報告である。
\dailyxsource{関連投稿}{https://x.com/CryptoMiners_Co/status/2122811995137413260}

\section{Vibe codingとコード理解を議論}
自然言語でアプリケーションを組み立てるvibe codingやClaude Codeを巡り、利用者の経験談と批判が交わされた。非エンジニアにも制作機会が広がる一方、生成されたコードを理解しないまま修正を重ねると、新たな不具合を招くとの懸念が示された。投稿ではClaude Codeの既存機能にも触れられているが、新規リリースとは区別が必要だ。Grokはトレンド規模や個別の体験談を独立確認していない。
\dailyxsource{コミュニティ投稿}{https://x.com/AGTPinsights/status/2122811896734618014}

\section{Sheetsでピクセルアートを生成}
Google DeepMindはGemini 4 HydrogenとAntigravityを用い、Google Sheetsの多数のセルを小さな画素に見立てたアニメーション付きピクセルアートのデモを紹介した。表計算ソフトを本来の集計用途から創作環境へ転用する例であり、エージェントが既存アプリケーションを操作する柔軟性を示す。数千セルを使ったとの説明があるが、具体的な再現手順や利用上限は観測窓内では示されていない。
\dailyxsource{Google DeepMind公式投稿}{https://x.com/GoogleDeepMind/status/2122807824531275826}